\section*{Capítulo \ref{ch:g10:drainage-basins} --- Bacias hidrográficas, rios e cheias}

\begin{solution}{exo:g10:drainage-basins:1}
A \omterm{def:g10:drainage-basins:drainage-basin}{bacia hidrográfica} de um ponto de um rio é toda a área de terra cujo
escoamento corre para o rio a montante desse ponto e passa por ele. O
\omterm{def:g10:drainage-basins:drainage-basin}{divisor de águas} é a linha, ao longo das cristas, que separa duas bacias
vizinhas. Duas gotas que caem a um metro uma da outra, de cada lado de um
divisor, descem por encostas opostas, para duas bacias diferentes: os seus
rios podem chegar a mares diferentes.
\end{solution}

\begin{solution}{exo:g10:drainage-basins:2}
A \omterm{def:g10:drainage-basins:discharge}{vazão} é o volume de água que passa por uma seção transversal do rio a
cada segundo, em \unit{m^3/s}. \qty{1}{m^3/s} são \qty{1000}{L} a cada
segundo; em um dia, $1000 \times 86\,400 = 86.4$ milhões de litros.
\end{solution}

\begin{solution}{exo:g10:drainage-basins:3}
Antes da cheia, o rio transportava cerca de \qty{330}{m^3/s} (entre cerca de
\num{240} e \qty{390}{m^3/s}). O pico foi de \qty{9231}{m^3/s} em 21 de
janeiro. A subida começou em 18 de janeiro (\qty{362}{m^3/s}): ela levou
três dias.
\end{solution}

\begin{solution}{exo:g10:drainage-basins:4}
$1/50 = 0.02$, \qty{2}{\percent}. Não: os anos são independentes; a chance
é de novo de \qty{2}{\percent} neste ano, aconteça o que tiver acontecido no
ano passado.
\end{solution}

\begin{solution}{exo:g10:drainage-basins:5}
Um pé cúbico vale $0.3048^3 \approx \qty{0.02832}{m^3}$, logo
$593\,000 \times 0.02832 \approx \qty{16800}{m^3/s}$. O Amazonas transporta
$206\,000/16\,800 \approx 12$ vezes mais.
\end{solution}

\begin{solution}{exo:g10:drainage-basins:6}
Áreas das faixas: $1.5 \times 0.3 = 0.45$, $0.90$, $1.05$ e \qty{0.60}{m^2}:
$A = \qty{3.0}{m^2}$. \omterm{def:g10:drainage-basins:discharge}{Vazões}: $0.2 \times 0.45 = 0.09$, $0.5 \times 0.90
= 0.45$, $0.6 \times 1.05 = 0.63$, $0.3 \times 0.60 = 0.18$: $Q =
\qty{1.35}{m^3/s}$. Velocidade média $1.35/3.0 = \qty{0.45}{m/s}$.
\end{solution}

\begin{solution}{exo:g10:drainage-basins:7}
$66 \times 10^6 \times 0.020 = \qty{1.32e6}{m^3}$. Dois dias são
$2 \times 86\,400 = \qty{172800}{s}$: $1.32 \times 10^6/172\,800 \approx
\qty{7.6}{m^3/s}$ em média.
\end{solution}

\begin{solution}{exo:g10:drainage-basins:8}
Ordenados: 120, 88, 65, 59, 51, 47, 42, 38, 35. Com $n = 9$, $T = 10/m$: o
maior (120) tem $T = 10$ anos, o segundo (88), 5 anos, o terceiro (65),
cerca de 3.3 anos. O mediano, 51, tem posto 5 e $T = 2$ anos: ele é igualado
ou superado um ano em cada dois.
\end{solution}

\begin{solution}{exo:g10:drainage-basins:9}
O \omterm{def:g10:drainage-basins:hydrograph}{hidrograma} da cidade sobe mais cedo e muito mais alto, e cai depressa: nas
encostas íngremes, nos telhados e nas ruas, pouca água se infiltra e o
escoamento chega ao córrego depressa e de uma vez; o seu \omterm{def:g10:drainage-basins:hydrograph}{tempo de retardo} é
curto. O \omterm{def:g10:drainage-basins:hydrograph}{hidrograma} da bacia florestada sobe mais tarde, até um pico mais
baixo e mais largo: o solo e a floresta absorvem boa parte da chuva e a
liberam devagar, e a água corre devagar nas encostas suaves. O seu
\omterm{def:g10:drainage-basins:hydrograph}{escoamento de base} fica mais alto depois da tempestade.
\end{solution}

\begin{solution}{exo:g10:drainage-basins:10}
$9231 \times 86\,400 \approx \qty{7.98e8}{m^3}$, espalhados sobre
$29\,900 \times 10^6 = \qty{2.99e10}{m^2}$: cerca de \qty{0.027}{m},
\qty{27}{mm}. O \omterm{def:g10:drainage-basins:hydrograph}{escoamento de base} dá $330 \times 86\,400/2.99 \times 10^{10}
\approx \qty{0.001}{m}$, cerca de \qty{1}{mm} por dia: o dia do pico
transportou cerca de 28 vezes mais.
\end{solution}

\begin{solution}{exo:g10:drainage-basins:11}
$P = 1 - 0.9^{10} \approx 1 - 0.349 = 0.651$: cerca de \qty{65}{\percent}.
Ela não vem nenhuma vez com probabilidade $0.9^{10} \approx 0.35$, cerca de
\qty{35}{\percent}.
\end{solution}

\begin{solution}{exo:g10:drainage-basins:12}
O \omterm{def:g10:drainage-basins:hydrograph}{hidrograma} é um triângulo de três dias de duração e \qty{1200}{m^3/s} de
altura: a sua área é $\tfrac12 \times 1200 \times 3 \times 86\,400 \approx
\qty{1.56e8}{m^3}$, \num{156} milhões de metros cúbicos. A parte acima de
\qty{800}{m^3/s} é um triângulo menor da mesma forma, três vezes menor em
altura ($400$ em vez de $1200$) e em duração (um dia em vez de três):
$\tfrac12 \times 400 \times 86\,400 \approx \qty{1.73e7}{m^3}$, \num{17}
milhões de metros cúbicos. Isso é menos que os \num{60} milhões que o
reservatório pode armazenar: sim.
\end{solution}

\begin{solution}{exo:g10:drainage-basins:13}
Cheio até as margens, o rio tem $800/(100 \times 2) = \qty{4}{m}$ de
profundidade. (1) A planície de inundação transporta \qty{400}{m^3/s} sobre
\qty{900}{m} a \qty{0.5}{m/s}: a sua profundidade é
$400/(900 \times 0.5) \approx \qty{0.9}{m}$. (2) Entre os diques:
$1200/(100 \times 2) = \qty{6}{m}$ de profundidade, \qty{2}{m} acima das suas margens, mais que o
dobro da profundidade na planície de inundação sem diques.
\end{solution}

\begin{solution}{exo:g10:drainage-basins:14}
Queremos $(1 - 1/T)^{30} = 0.5$, logo $1 - 1/T = 0.5^{1/30} \approx 0.97716$,
$1/T \approx 0.02284$ e $T \approx 44$ anos.
\end{solution}

\begin{solution}{exo:g10:drainage-basins:15}
Cada quadrado manda a sua água para o seu vizinho mais baixo (linhas de
norte a sul, colunas de oeste a leste):
\begin{center}
\begin{tabular}{|c|c|c|c|}
\hline
exutório & oeste & oeste & oeste \\ \hline
norte & oeste & oeste & sul \\ \hline
norte & norte & sul & sul \\ \hline
norte & leste & leste & exutório \\ \hline
\end{tabular}
\end{center}
Seguindo as setas, os quadrados 47, 39, 44, 37, 46 e 30 drenam para o
exutório sudeste: 6 quadrados, \qty{0.06}{km^2}. Os outros 10 quadrados
drenam para o exutório noroeste: \qty{0.10}{km^2}. O divisor vai da borda
leste, entre 50 e 47, depois entre 48 e 47, entre 48 e 44, entre 49 e 44,
entre 49 e 46, e entre 51 e 46, até a borda sul.
\end{solution}

\begin{solution}{pb:g10:drainage-basins:1}
\textbf{1.} $11\,560 \times 1.609^2 \approx \qty{29900}{km^2}$.

\textbf{2.} O \omterm{def:g10:drainage-basins:drainage-basin}{divisor de águas} da bacia a montante da estação.

\textbf{3.} $484\,000 \times 0.3048^3 \approx \qty{13700}{m^3/s}$.

\textbf{4.} $13\,705/3341 \approx 4.1$ vezes.

\textbf{5.} $A = 400 \times 8 = \qty{3200}{m^2}$ e $v = Q/A = 13\,705/3200
\approx \qty{4.3}{m/s}$.

\textbf{6.} Cerca de \qty{330}{m^3/s}.

\textbf{7.} \qty{9231}{m^3/s}, em 21 de janeiro.

\textbf{8.} Três dias.

\textbf{9.} $9231 \times 86\,400/(29\,900 \times 10^6) \approx
\qty{0.027}{m}$: cerca de \qty{27}{mm}.

\textbf{10.} Água nova chegou ao rio: outra chuva sobre a bacia (ou neve
derretendo) mandou um segundo pulso, menor, de escoamento.

\textbf{11.} Dezessete dias depois do pico. A subida vem do escoamento das
encostas perto dos córregos, que chega depressa e de uma vez; no ramo
descendente, o rio ainda recebe a água das partes mais distantes da bacia e
a água que escoa devagar do solo.

\textbf{12.} $T = n + 1 = 96$ anos.

\textbf{13.} $96/2 = 48$ anos.

\textbf{14.} $96/3 = 32$ anos.

\textbf{15.} $96/4 = 24$ anos e $96/5 = 19.2$ anos. Não muito: de um modo
ou de outro, essa \omterm{def:g10:drainage-basins:discharge}{vazão} volta cerca de uma vez a cada vinte anos.

\textbf{16.} Cinco anos (1936, 1937, 1943, 1972 e 1996): $5/95 \approx
0.053$, cerca de \qty{5}{\percent} por ano, um \omterm{def:g10:drainage-basins:return-period}{período de retorno} de cerca de
$96/5 \approx 19$ anos.

\textbf{17.} $1/100 = 0.01$, \qty{1}{\percent}.

\textbf{18.} $1 - 0.01 = 0.99$.

\textbf{19.} $0.99^{30} \approx 0.740$.

\textbf{20.} $1 - 0.740 = 0.26$: \textbf{a chance de pelo menos uma cheia de
100 anos em 30 anos é de cerca de \qty{26}{\percent}}, uma em quatro.
\end{solution}
